% !TEX root = ../nc-cubic.tex
\subsection{任意個の中心変数に対する二捩れ型の比較}
\label{nsf:sec:general}
\begin{theorem}\label{nsf:thm:all-dimensions}
$m\ge2$、$U=(u_1,\ldots,u_m)$ とし、$p,q,l$ は $U$ の非零一次式、
$g$ は $U$ の三次式、$\lambda\ne0$ とする。$G,\alpha$ は前節のものとし、
$G$ は $U$ を固定する。$F=px^2+qy^2+lz^2+\lambda xyz+g$ に対して
\[
 F'=px'^2+qy'^2-lz'^2+\lambda x'y'z'+g+4pql/\lambda^2
\]
と置く。$V(F)\subset\PP^{m+2}$ が滑らかなら $V(F')$ も滑らかであり、
\[
 \Ku(V(F))^{G,\alpha}\simeq\Ku(V(F'))
\]
という $\C$-線型 dg Morita 同値がある。左辺は三つの反可換変数と
$m$ 個の中心変数を持つ対応する量子中心三次商の全 coherent 残余圏である。
\end{theorem}
\begin{proof}
滑らかさの計算は底の次元に依らず次の形になる。
$P=pql$, $h=g+4P/\lambda^2$ とする。
$U\ne0$ で active 座標が零個または三個の特異条件は、
それぞれ $dg=0$ または $P\ne0,\ dh=0$ である。
一個だけの場合を零個の場合と合わせれば、いずれかの $L=p,q,l$ に対して
$g=L=0$ かつ $dg$ が $dL$ のスカラー倍となることも排除すべきである。
二個だけの場合と $U=0$ は前定理と同じである。
$L=0$ 上で $h=g$、$dh-dg$ は $dL$ のスカラー倍なので、
全条件は $(g,h,l)\mapsto(h,g,-l)$ で不変となる。

Orlov の整数次数付き行列因子化模型を使う。
同変側の同値でも、次数による切り詰め、自由加群による射影および
線束の例外部分は $G$-線型作用を保つ。有限群の同変化は
中心拡大の指定文字成分として取る。
全 coherent 圏について量子表示との同定は定理\ref{nqt:thm:reduction}の
行列不変環と全次数シフトの射影生成族の証明による。
非退化な乗数により各例外部分で既約射影表現は一つである。

$r=\lambda z/2$, $s=\lambda z'/2$ とし、共通の係数環
$H=\C[U,d]$ において $\deg d=2$ とする。
二つの多項式底 $R=\C[U,r]$, $R'=\C[U,s]$ は
$d\mapsto r^2$, $d\mapsto pq-s^2$ でそれぞれ $H$ 上自由階数二である。
$C=g+4ld/\lambda^2$ を共通の曲率とする。残る二変数の二次部分は
$Q=px^2+2rxy+qy^2$ および $Q'=px'^2+2sx'y'+qy'^2$ である。

\emph{有限生成対象。}
二重次数（コホモロジー次数、内部次数）で周期元 $t,t^{-1}$ を用い、
$|t|=(2,-3)$、曲率は $tF$ とする。これは二つの自由加群を
$(3)=[2]$ により並べた通常の graded MF である。
$F=\sum_j A_jB_j$ の Koszul 因子化の微分を
$\sum_j(A_j\iota_j+t B_j\epsilon_j)$ と書く。
$\epsilon_j$ は外積で次数 $(-1,\deg A_j)$、$\iota_j$ はその収縮である。
$A_j$ が周囲環で正則列なら、この因子化はその列の商加群の
stabilization である。実際、通常の Koszul 解消に $\sum B_j\epsilon_j$
という平方零の乗法ホモトピーを入れた Shamash 解消を作ると、
十分先の二周期部分がこの微分になる。

$K_C$ は $H$ 上で $(U,d)$ の Koszul 因子化とする。
$g=\sum u_i g_i$ と書き最後の係数を $4l/\lambda^2$ とすればよい。
$K_Q$ は $R[x,y]$ 上で $(x,y)$ と $(px+ry,rx+qy)$ を用いるものとする。
$K=K_Q\otimes_H K_C$（自然な周期元を共通に取る）は $F$ の因子化で、
$G$ は $K_C$ を固定する。これは $(x,y,U,r^2)$ の商
$M=\C[r]/(r^2)$ の stabilization である。
双対側にも $K'=K_{Q'}\otimes_H K_C$ があり、
$(x',y',U,pq-s^2)$ の商 $M'=\C[s]/(s^2)$ を与える。

$K'$ の内部シフト全体は双対側の全 MF を厚生成する。
実際、$F'$ は $s$ 軸全体で消える。この軸の環 $\C[s]$ の次数付き
導来圏で $s:M'(-1)\to M'$ の cone は、その二つのコホモロジーの
shift の直和である。これは $\C[s]$ の大域次元が一であるためで、
実際、核は $sM'(-1)=\C(-2)$、余核は $\C$ なので
cone は $\C\oplus\C(-2)[1]$ である。軸からの加群の制限と
特異圏への射影はこの分裂を保つ。
滑らかな射影超曲面では剰余体の内部シフトが全 graded 特異圏を生成する。
実際、Jacobian イデアルは各変数のある冪を含み、微分行列を微分すると
その偏導関数の乗法は MF 上 null-homotopic である。
変数の個数を $N=m+3$ とすると勾配は二次の正則列であり、
Jacobian 環の Hilbert 級数は $(1+z)^N$ だから、各冪は一様に $N+1$ でよい。
各変数のその冪による Koszul 複体を対象に tensor すれば元の対象は
そのシフトを含む直和の因子になる。一方、この複体で切った対象は
特異圏に移す前に原点支持の有限長コホモロジーを持ち、剰余体から生成される。
源側でも同じ議論で、作用を忘れた $K$ のシフトが生成する。
ねじれ誘導・忘却と平均化により $K\otimes V$ のシフトは
$\alpha$-同変全圏を生成する。ここで $K$ 自身は通常の $G$-線型化を持ち、
ねじれ誘導は $K$ とねじれ正則表現 $V^{\oplus2}$ の tensor である。
具体的には任意のねじれ同変対象はその忘却のねじれ誘導の直和因子となる。
余単位の各群元成分に線型化の逆を掛けて $1/|G|$ で平均すると切断を得る。
従って $(3)=[2]$ も用いれば、源では $K\otimes V$、標的では $K'$ の
内部シフト $0,1,2$ の直和一個ずつで全圏を厚生成する。

\emph{二変数の微分付き自己準同型。}
$Q_{ij}$ を行列 $\left(\begin{smallmatrix}p&r\\r&q\end{smallmatrix}\right)$
の成分とし、$\theta_i=\iota_i-t\sum_j Q_{ij}\epsilon_j$ と置く。
$\operatorname{End}(K_Q)$ の超交換子微分では
$\partial\epsilon_1=x,\ \partial\epsilon_2=y,\ \partial\theta_i=0$ である。
$\theta_i$ の生成する Clifford 代数 $D_r$ は係数 $R[t^{\pm1}]$ 上
反交換子 $-2tQ_{ij}$ を持つ。零微分の $D_r$ からの包含
\[
             D_r\longrightarrow\operatorname{End}(K_Q)
\]
は二重次数付き dg 準同型である。
実際、右辺は $R[x,y,t^{\pm1}]$ 上の基底
$\epsilon^I\theta^J$ を持つ。これは通常の $\epsilon^I\iota^J$ から
収縮の個数による三角基底変更で得られる。
係数複体は各 $J$ で $(x,y)$ の Koszul 複体なので、その
$R[t^{\pm1}]$ への収縮は $R[t^{\pm1}]$-線型に取れる。
一変数では $x^i\mapsto x^{i-1}\epsilon$ ($i>0$) でよく、二つを tensor する。
従って $\operatorname{End}_H(K_C)$ を周期元を共有して tensor した後も
包含は準同型である。全自己準同型がこの tensor 積となるのは、
両因子が有限自由で外積の微分が二つの微分の和だからである。
この証明は内部シフト間の全 Hom を同時に保つ。
$K\otimes V$ の同変自己準同型ではさらに $\operatorname{End}V$ を
tensor して不変部分を取る。有限群の平均化によりこの操作は完全である。

\emph{乗数を含む Clifford 同型。}
$\tau=-t$ とし、$a=\theta_1\theta_2/\tau-r$ と置けば $a^2=r^2-pq$。
前節の $J,K,M$ を用いる（ここでは $K$ の文字を行列に使う）。
$(D_r\otimes\operatorname{End}V)^G$ の中で中心的な $s=aM$ は
$s^2=pq-r^2$ を満たす。
$\theta'_1=\theta_1J,\ \theta'_2=\theta_2K$ は
平方 $\tau p,\tau q$、反交換子 $2\tau s$ を持つ。
従って双対二変数の Clifford 代数 $D_s$ から
$(D_r\otimes\operatorname{End}V)^G$ への射がある。
これは $H[t^{\pm1}]$ 上の次数を保つ同型である。
実際、$a'=\theta'_1\theta'_2/\tau-s=rM$ とすると、
前節の基底変更と同じく $1,a',\theta'_1,\theta'_2$ と
その $s$ 倍は不変部分の八項基底になる。
具体的には $sa'=-ra$, $s\theta'_1=-r\theta_1 K+p\theta'_2$、
$s\theta'_2=q\theta'_1-r\theta_2J$ である。
$d$ の像は両側で同じであり、$\operatorname{End}_H(K_C)$ の底変更にも
この同型を適用できる。ここで $G$ はこの最後の因子に自明に作用し、
その各項は $H[t^{\pm1}]$ 上自由なので不変部分と tensor 積は可換する。
以上の同型と準同型で、両側の生成対象の内部シフト間の
すべての Hom と積が対応する。これらの dg 部分圏の厚生成・冪等完備化により
主張の実際の dg Morita 同値が得られる。
\end{proof}
